\chapter{Experimental History}\label{ch:history}

This chapter tells the research story in order. Each phase is described by what motivated it, what was done, what was
found, and what it led to. Dates come from the dated entries of \fpath{docs/TASKS.md}, \fpath{docs/DECISIONS.md} and
\fpath{docs/RESULTS.md}. Work before 2026-11-04 (the original blueprint, environment and data set-up) was recorded only
as completed tasks, without dates. The mid-semester presentation that proposed the first extension ideas is dated
2026-10-12.

\begin{figure}[htbp]
\centering
\begin{adjustbox}{max width=\linewidth}\begin{tikzpicture}[x=1.13cm,y=1cm,
  ev/.style={font=\sffamily\scriptsize,align=center,text width=1.95cm},
  bar/.style={rounded corners=1.5pt,minimum height=0.32cm,inner sep=0pt},
  lab/.style={font=\sffamily\tiny,text=white,anchor=west}]
  % axis: day 4 .. 17 of November 2026
  \draw[thick,navy,-{Stealth[length=5pt]}] (3.7,0) -- (17.4,0);
  \foreach \d in {4,...,17} {\draw[navy] (\d,0.08) -- (\d,-0.08); \node[font=\sffamily\tiny,below=2pt] at (\d,0) {\d};}
  \node[font=\sffamily\scriptsize,anchor=west] at (3.7,-0.55) {November 2026 (IST)};
  % decisions and results (upper lane)
  \foreach \x/\y/\t in {4/1.1/{Gate fails on base; reversal to LoRA, Scala, 3 rounds},
                         7/2.75/{Core study done: TS beats NV at round 3},
                         8/1.1/{DPO harms; repair pilot: Instruct repairer go},
                         10/2.75/{Repair-TS base 0.5856; plan A/B approved},
                         12/1.1/{Instruct TS matches paper},
                         13/2.75/{Controls: repair null on Instruct SQL},
                         15/1.1/{Repair-ED built},
                         16/2.75/{Scala done: repair wins every round; NL-to-R built},
                         17/1.1/{ED and NL-to-R done: repair-TS-R wins every round}}{
    \draw[sea,thick] (\x,0) -- (\x,\y-0.25);
    \fill[sea] (\x,0) circle (2pt);
    \node[ev,anchor=south] at (\x,\y-0.3) {\t};}
  % experiment bars (lower lane)
  \node[bar,fill=navy,minimum width=3.6cm,anchor=west] at (4.95,-1.2) {}; \node[lab] at (5.0,-1.2) {Core study TS/NV, base};
  \node[bar,fill=plum,minimum width=1.5cm,anchor=west] at (7.85,-1.65) {}; \node[lab] at (7.9,-1.65) {Ext.~2};
  \node[bar,fill=sea,minimum width=2.6cm,anchor=west] at (8.4,-2.1) {}; \node[lab] at (8.45,-2.1) {Repair-TS base};
  \node[bar,fill=teal,minimum width=2.6cm,anchor=west] at (10.9,-1.2) {}; \node[lab] at (10.95,-1.2) {Instruct TS/NV};
  \node[bar,fill=amber,minimum width=2.4cm,anchor=west] at (12.75,-1.65) {}; \node[lab,text=black] at (12.8,-1.65) {Repair-TS Instr.};
  \node[bar,fill=coral!80,minimum width=2.1cm,anchor=west] at (14.7,-1.2) {}; \node[lab] at (14.75,-1.2) {Scala};
  \node[bar,fill=slate,minimum width=3.0cm,anchor=west] at (14.75,-2.1) {}; \node[lab] at (14.8,-2.1) {ED TS, repair-ED};
  \node[bar,fill=plum,minimum width=1.7cm,anchor=west] at (16.2,-1.65) {}; \node[lab] at (16.25,-1.65) {NL-to-R};
  % incidents
  \foreach \x/\t in {5.4/crash, 14.8/stall, 15.5/WSL} {
    \fill[coral] (\x,0) circle (2.5pt);
    \node[font=\sffamily\tiny,text=coral,above=1pt] at (\x,0.05) {\t};}
\end{tikzpicture}
\end{adjustbox}
\caption{Project timeline, 2026-11-04 to 2026-11-17 (IST). Upper lane: decisions and results; lower lane: experiments
and their durations; red markers: machine incidents.}
\label{fig:timeline}
\end{figure}

\section{Phase 0: the blueprint and its reversal (before and on 2026-11-04)}

\begin{description}[style=nextline,leftmargin=1.2em]
\item[Plan.] The first plan, written under a deadline, was deliberately reduced in scope:
  \begin{itemize}
  \item Llama-3.1-8B-Instruct trained with QLoRA (4-bit weights), estimated to fit 24~GB with margin;
  \item Pandas as the second language, because the repository has an NL-to-Pandas generator and a fast Python sandbox;
  \item two rounds;
  \item error detection and ablations only if time remained.
  \end{itemize}
\item[Reversal (2026-11-04).] The deadline was removed, so every choice made to save time or memory was reopened in
favour of faithfulness:
  \begin{itemize}
  \item \textbf{LoRA} as in the paper, measured at 20.0~GB at about 4,000 tokens: it fits.
  \item \textbf{Spark-Scala} as in the paper, at a measured cost of about 50 times Pandas' run time per program.
  \item \textbf{Three rounds.}
  \item \textbf{Base checkpoint}, because the README links it (\cref{sec:design}).
  \end{itemize}
\item[Environment results.]
  \begin{itemize}
  \item LoRA on base peaked at 19.8~GB at 4,000 tokens.
  \item vLLM served base at 7,417 prefill and 262 decode tokens/s.
  \item The Spark sandbox first ran 0 of 32 programs: Spark's JVM hit the sandbox's 256 open-file limit. With the limit
  raised to 4,096, it ran 48 of 48.
  \item Spark throughput was 2,133 programs/h at 24 workers and 2,714 at 36.
  \end{itemize}
\end{description}

\section{Phase 1: the baseline gate fails on base (2026-11-04)}\label{sec:hist-gate}

\begin{description}[style=nextline,leftmargin=1.2em]
\item[Motivation.] Before any fine-tuning, check that vanilla scores are within $\pm0.05$ of the paper's. Otherwise
fine-tuned gains cannot be compared.
\item[Result.] Base with raw completions and the zero-shot prompt met the gate only on BIRD (0.3636 against 0.388).
  \begin{itemize}
  \item WikiSQL: 0.326 against 0.525. 507 of 1,000 answers were SQL that fails to run, often from unquoted multi-word
  column names.
  \item WikiTQ: 0.219 against 0.300.
  \item Spider: 0.7245 against 0.670, too high.
  \end{itemize}
Six setting combinations were tried in all (\cref{sec:fail-gate}); none reproduced the paper.
\item[Decision.] Stop searching and proceed with the gate failed, recording the gaps as unexplained. Further
combinations would eventually land near 0.525 by chance. The report would compare the ordering of the arms, not
absolute scores. The authors could not be contacted, so the questions for them stay open (\cref{sec:open-questions}).
\end{description}

\section{Phase 2: building and auditing the validator (2026-11-04)}\label{sec:hist-validator}

\begin{description}[style=nextline,leftmargin=1.2em]
\item[Motivation.] The repository's validation never compares SQL with Scala results, so the comparator is ours, and a
silent bug in it would decide every training set.
\item[Action.] A 300-table, 2-round test run, followed by a hand inspection of 50 kept and 24 rejected pairs. All 160
\texttt{scala\_failed} pairs were re-run for their error text.
\item[Findings.] Five comparator defects (\cref{sec:comparator}). After the fixes, a second 300-table run and a fresh
inspection of 50 kept and 50 rejected pairs found:
  \begin{itemize}
  \item about 8 in 10 kept pairs fit their question;
  \item about 1 in 3 rejected pairs may have held usable SQL.
  \end{itemize}
\item[Consequence.] The core study ran on the corrected validator. The inspection's second finding later motivated
both extensions.
\end{description}

\section{Phase 3: the core study on base (2026-11-04 to 2026-11-07)}

\begin{description}[style=nextline,leftmargin=1.2em]
\item[Action.] TS and NV, 3 seeds $\times$ 3 rounds, one arm per GPU, started 2026-11-04 at about 23:25.
\item[Incidents.] On 2026-11-05, a U+2028 character in a training file stopped NV, and the machine crashed while both
arms were training. Both runs resumed from checkpoints. These incidents produced the project's standing operating
rules: \texttt{tmux}, append-only logs, one job per GPU, checkpoint copies, and no commits (\cref{app:incidents}).
\item[Result.] Both arms finished on 2026-11-07. TS beat NV on all 8 round-3 cells; NV beat TS in round~1 in 17 of 24
seed--cell comparisons (\cref{sec:sql-core}).
\item[Consequence.] The paper's qualitative claim was reproduced on base. The user chose Extension~2 as the first
contribution beyond the paper.
\end{description}

\section{Phase 4: Extension 2 fails, Extension 3 succeeds (2026-11-07 to 2026-11-10)}

\begin{description}[style=nextline,leftmargin=1.2em]
\item[Extension 2 (bestof4, DPO).]
  \begin{itemize}
  \item Sampling 3 extra candidates per question found a validated program for 52--62\% more questions, yet moved the
  test mean by only $+0.0064$.
  \item DPO lowered it by $-0.0209$.
  \item Tuning five gentler DPO variants on dev sets under a pre-set rule could not beat bestof4 (\cref{sec:ext2}).
  \item The remaining rounds of Extension~2 were put on hold.
  \end{itemize}
\item[Extension 3 (repair).] The user chose repair from a list of proposed ideas.
  \begin{itemize}
  \item \textbf{Base repairer pilot}: +17.1\% kept pairs, below the pre-set +30\% threshold, so no-go. Base copied the
  failed code back in half of Scala repairs.
  \item \textbf{Instruct repairer pilot}: +56.1\%, with 21 of 23 recovered pairs correct by hand, so go.
  \item \textbf{Full round~1}: beat TS round~1 on every seed.
  \item \textbf{Rounds 2--3}: reached a mean of 8 of 0.5856 against TS's 0.5435.
  \end{itemize}
\end{description}

\section{Phase 5: matching the paper with Instruct (2026-11-10 to 2026-11-12)}

\begin{description}[style=nextline,leftmargin=1.2em]
\item[Motivation.] A plan analysing the paper, the repository and our own results (\fpath{docs/PLAN_NEXT.md}) observed
three things:
  \begin{itemize}
  \item the repository's only query path is chat with a fixed system message;
  \item earlier Instruct scores through that path were within 0.011 of the paper on WikiSQL and WikiTQ;
  \item base cannot use that path at all.
  \end{itemize}
The plan was approved on 2026-11-10.
\item[Result.] Instruct vanilla passed the gate on 3 of 4 sets, all within 0.015. Instruct TS (3 seeds $\times$ 3
rounds) and Instruct NV (seed~1) finished on 2026-11-12. TS reproduced the paper's Table~3 within 0.032 per set,
Spider excepted (\cref{sec:sql-instruct}).
\item[Alongside.] Self-consistency, an established test-time method, was tuned on dev and applied to every final model
(\cref{sec:ext-sc}). A BIRD scoring artefact was found and fixed (\cref{sec:fail-bird}).
\end{description}

\section{Phase 6: repair-TS on Instruct and the controls (2026-11-12 to 2026-11-14)}

\begin{description}[style=nextline,leftmargin=1.2em]
\item[Motivation.] Does repair-TS's base gain carry over to the model that matches the paper? And how much of it is
just more data?
\item[Result.]
  \begin{itemize}
  \item Instruct repair-TS showed no gain at any round (round~3: $-0.0014$, $p=0.55$).
  \item The same-data control split base repair-TS's round-1 gain about in half between volume and content.
  \item On Instruct, extra first-pass pairs helped while repaired pairs did not (\cref{sec:ext-control}).
  \end{itemize}
\item[Consequence.] Repair-TS was reframed from a method into a conditional finding.
\end{description}

\section{Phase 7: the paper's other tasks (2026-11-14 to 2026-11-17)}

\begin{description}[style=nextline,leftmargin=1.2em]
\item[Motivation.] Extend the reproduction to other tasks, ranked by feasibility. NL-to-Scala needed no new training.
Error detection had downloadable test data. NL-to-R needed new data generation and an R sandbox.
\item[NL-to-Scala.] 19 models $\times$ 8 cells were scored, finishing on 2026-11-16. Vanilla matched the paper's
scores on 3 of 4 sets. Repair-TS beat TS at every round (\cref{sec:results-scala}). A quadratic
comparison discovered while scoring BIRD was fixed with no change to any verdict (\cref{sec:fail-quadratic}).
\item[Error detection.] The vanilla gate passed within 0.006. TS ran 3 seeds $\times$ 3 rounds. Repair-ED (3 seeds $\times$
3 rounds) and its two round-1 controls ran from 2026-11-15 to 2026-11-17 (\cref{sec:results-ed}).
\item[NL-to-R.] Approved and built on 2026-11-16, smoke-tested the same day (which exposed truncated R error feedback,
fixed before the full runs), and finished on 2026-11-17: TS-R and repair-TS-R, 3 seeds $\times$ 3 rounds, 19 models
scored (\cref{sec:results-r}).
\item[Incidents.] A 7.5-hour machine stall on 2026-11-14 and a WSL crash on 2026-11-15. Both were recovered with no
lost result (\cref{app:incidents}).
\end{description}

\section{How each phase led to the next}

\begin{enumerate}
\item The \textbf{failed gate} (phase 1) made absolute comparisons with the paper unreliable on base. Phase 3
therefore focused on the paper's \emph{relative} claim (TS versus NV versus vanilla), and it later motivated trying
the repository's own query path on Instruct (phase 5).
\item The \textbf{validator audit} (phase 2) found that about a third of rejections might be recoverable. This
motivated both extensions (phase 4).
\item \textbf{DPO's failure} (phase 4) showed that penalising rejected programs drifts the model. It redirected effort
from learning \emph{against} rejected data to \emph{recovering} it.
\item \textbf{Base's failure as a repairer} (phase 4) led to the Instruct repairer. That introduced a second model
into the base track, and the distillation confound discussed in \cref{ch:discussion}.
\item \textbf{Instruct matching the paper} (phase 5) raised the question of whether repair-TS still helps on that
model (phase 6). It did not for SQL, but it did for Scala (phase 7).
\end{enumerate}
